\begin{figure}[H]
\centering
\definecolor{phasePipeline}{HTML}{2563EB}
\definecolor{phaseReference}{HTML}{059669}
\definecolor{phaseTotal}{HTML}{111827}
\definecolor{phaseQaoa}{HTML}{D97706}
\definecolor{phaseGrid}{HTML}{E5E7EB}
\definecolor{phaseText}{HTML}{374151}
\begin{tikzpicture}
\begin{axis}[
    width=0.90\textwidth,
    height=0.53\textwidth,
    xmode=log,
    ymode=log,
    xmin=9,
    xmax=1100,
    ymin=0.000003,
    ymax=5,
    xlabel={$N$},
    ylabel={Desequilíbrio normalizado (\%)},
    xlabel style={font=\small\sffamily, text=phaseText},
    ylabel style={font=\small\sffamily, text=phasePipeline},
    tick label style={font=\footnotesize\sffamily, text=phaseText},
    xtick={10,20,50,100,200,500,1000},
    xticklabels={10,20,50,100,200,500,1000},
    ytick={0.00001,0.0001,0.001,0.01,0.1,1},
    yticklabels={0.00001,0.0001,0.001,0.01,0.1,1},
    log basis x=10,
    log basis y=10,
    grid=both,
    major grid style={draw=phaseGrid, line width=0.35pt},
    minor grid style={draw=phaseGrid!45, line width=0.25pt},
    axis line style={draw=phaseText!35},
    tick style={draw=phaseText!45},
    axis y line*=left,
    axis x line*=bottom,
    legend style={
      draw=none,
      fill=none,
      font=\footnotesize\sffamily,
      at={(0.02,0.98)},
      anchor=north west,
      cells={anchor=west}
    },
    legend image post style={line width=1.2pt},
]
\addplot[
    color=phasePipeline,
    line width=1.15pt,
    mark=*,
    mark size=1.8pt,
    mark options={fill=white, draw=phasePipeline, line width=0.8pt}
] table[x=N,y=pipeline_quality_pct,col sep=comma]
  {data/phase5_scalability_plot.csv};
\addlegendentry{Pipeline}

\addplot[
    color=phaseReference,
    line width=1.05pt,
    dashed,
    mark=diamond*,
    mark size=1.8pt,
    mark options={fill=white, draw=phaseReference, line width=0.75pt}
] table[x=N,y=reference_quality_pct,col sep=comma]
  {data/phase5_scalability_plot.csv};
\addlegendentry{Melhor referência}
\end{axis}

\begin{axis}[
    width=0.90\textwidth,
    height=0.53\textwidth,
    xmode=log,
    ymode=log,
    xmin=9,
    xmax=1100,
    ymin=3,
    ymax=900,
    ylabel={Tempo (s)},
    ylabel style={font=\small\sffamily, text=phaseTotal},
    tick label style={font=\footnotesize\sffamily, text=phaseText},
    xtick=\empty,
    ytick={5,10,30,100,300,900},
    yticklabels={5,10,30,100,300,900},
    log basis x=10,
    log basis y=10,
    axis y line*=right,
    axis x line=none,
    axis line style={draw=phaseText!35},
    tick style={draw=phaseText!45},
    legend style={
      draw=none,
      fill=none,
      font=\footnotesize\sffamily,
      at={(0.98,0.98)},
      anchor=north east,
      cells={anchor=west}
    },
    legend image post style={line width=1.2pt},
]
\addplot[
    color=phaseTotal,
    line width=1.15pt,
    mark=square*,
    mark size=1.7pt,
    mark options={fill=white, draw=phaseTotal, line width=0.8pt}
] table[x=N,y=total_time_s,col sep=comma]
  {data/phase5_scalability_plot.csv};
\addlegendentry{Pipeline total}

\addplot[
    color=phaseQaoa,
    line width=1.15pt,
    mark=triangle*,
    mark size=2.0pt,
    mark options={fill=white, draw=phaseQaoa, line width=0.8pt}
] table[x=N,y=qaoa_time_s,col sep=comma]
  {data/phase5_scalability_plot.csv};
\addlegendentry{QAOA}

\node[
    anchor=south east,
    align=right,
    font=\scriptsize\sffamily,
    text=phaseText!78,
    fill=white,
    fill opacity=0.88,
    text opacity=1,
    inner xsep=5pt,
    inner ysep=4pt
] at (rel axis cs:0.985,0.055)
{XY, $p=2$, \texttt{top\_k}=10\\
\texttt{qubit\_max}=8; sub-QUBOs: 3--250\\
BF até $N=20$; depois SA/MILP\\
viabilidade global: 100\%};
\end{axis}
\end{tikzpicture}
\caption{Escalabilidade corrigida da pipeline iterativa na Fase 5. O eixo esquerdo compara o desequilíbrio normalizado médio da pipeline com a melhor referência disponível: brute force nas instâncias pequenas e simulated annealing ou testemunho MILP nas restantes. O eixo direito mostra o tempo médio total e o tempo médio consumido pelo QAOA, ambos em escala logarítmica.}
\label{fig:phase5-scalability-dual-axis}
\end{figure}
